% 13-operations.tex: выживание в эксплуатации.
\section{Выживание в действительности: стек, отказы, нагрузка, постоянный свидетель и день, когда заканчивается пилот}
\label{sec:operations}

\subsection{Стек}

Промышленная топология состоит из пяти служб: собственной сборки кромки ТЛС, которая согласует с
современными клиентами гибридный постквантовый обмен ключами; самого приложения под управлением
сервера гуникорн; пулера соединений; ПостгреСКЛ с непрерывным архивированием; и Редиса для общего
ограничителя частоты запросов. Каждая служба работает от непривилегированного пользователя со
снятыми возможностями ядра, сторонние образы закреплены по дайджесту, и ни один промышленный
контейнер не несёт в себе средств тестирования. Четыре пути развёртывания разделяют одну
дисциплину секретов, один путь резервного копирования и одну модель угроз: запуск на ноутбуке;
профиль для одного хоста с поочерёдным <<сине-зелёным>> выпуском; сценарная установка на Линукс под
управлением системд, проверяемая на двух дистрибутивах; и эталонный профиль Хелм для Кубернетеса
с сетевыми политиками <<по умолчанию всё запрещено>> и строгим стандартом безопасности подов,
поднимаемый в НИ на кластере из одного узла. \sImpl

\subsection{Отказ, отработанный и измеренный}

Правило репозитория гласит: проверять запуском, а не чтением. Из десяти дефектов, найденных за
один обход, чтением не был бы замечен ни один. Каждое эксплуатационное утверждение ниже есть учение,
которое выполняется в НИ и фиксирует число.

\begin{description}
  \item[Отказоустойчивость.] База данных работает под Патрони, с арендой ведущего в ЕТЦД из трёх
    участников и маршрутизацией по ролям. Под живым потоком записей учение теряет ведущий узел
    (реплика повышается по истечении двадцатисекундной аренды, выполняемые запросы быстро падают и
    повторяются, ни один не потерян), отрезает ведущий узел от хранилища аренды (он уступил за
    7~секунд), выполняет плановое переключение (3,3~секунды, ни одна вставка не упала) и роняет
    участника ЕТЦД (остановка на 0,3~секунды). Профиль Кубернетеса запускает тех же участников,
    используя в качестве хранилища аренды программный интерфейс самого кластера, и удаляет под
    ведущего узла. \sImpl
  \item[Второй регион.] Профиль отказоустойчивости переживает гибель узла, но не гибель региона, и
    соблазнительное решение, третий участник в регионе Б, неверно сразу по трём причинам: оно
    помещает глобальную сеть внутрь кворума хранилища аренды, на путь каждой фиксации и в зону
    поражения региона А. Поэтому начиная с \code{в9.359} регион Б есть резервный кластер: со своим
    хранилищем аренды и своей областью, получающий поток изменений асинхронно через маршрутизатор
    региона А, так что переключение внутри А не рвёт репликацию, и не выбирающий собственного
    ведущего, так что два региона никогда не принимают записи одновременно. Точка восстановления
    поэтому не нулевая, и она измеряется, а не провозглашается: учение гасит регион А так, как гаснет
    настоящий регион, повышает Б под живым потоком записей и сообщает время до начала обслуживания
    (5~секунд) и число строк, подтверждённых в А и отсутствующих в Б (ноль при темпе учения), против
    потолков в 90~секунд и 50~строк. Оно утверждает и две вещи поважнее любого числа: в Б нет ни
    одной строки, которую А не подтверждал, и дошедшее есть непрерывное начало последовательности,
    потому что резерв, держащий строки с первой по сороковую и с сорок пятой по шестидесятую, не
    отстал, а пропустил. Вернувшийся регион А перестраивается как резервный и никогда не
    возвращается ведущим. \sImpl
  \item[Реплики для чтения.] Поверхности, которые только читают, направляются на потоковую реплику
    по явному договору об отставании, по умолчанию не более десяти секунд, а сверх того запрос
    уходит на ведущий узел; каждый такой ответ называет свой источник и отставание. Чтения, от
    которых зависит правильность, остаются на ведущем узле, и половина вердикта проверки, отвечающая
    за полномочия, закреплена именно там (раздел~\ref{sec:authorization}). \sImpl
  \item[Разбиение на разделы.] Четыре таблицы событий, работающие только на добавление, разбиты на
    помесячные разделы, и учение доказывает, что триггер добавления держит и в разделе, и при
    присоединении, и при отсоединении, потому что исключение, сделанное для архивов, не должно
    стать путём к переписыванию. \sImpl
  \item[Аварийное восстановление.] Резервные копии зашифрованы, и восстановление без ключа
    падает закрытым. Ежемесячное учение убивает ведущий узел, восстанавливает его из архива и
    фиксирует измеренные числа против целей в 300~секунд потерянных данных и четыре часа на
    восстановление: последние записанные строки показывают от 36 до 42~секунд потерь и меньше пяти
    секунд до начала обслуживания, на одноразовом хосте учения, а не на промышленном железе. \sImpl
  \item[Хаос.] Еженедельное учение поднимает сине-зелёный стек под трафиком, убивает один цвет без
    единого потерянного запроса, останавливает оба, пока страница о сбое не дойдёт до настоящего
    диспетчера оповещений и его обратного вызова, убивает Редис и ПостгреСКЛ, отрезает пулер от сети
    и фиксирует каждое время восстановления против потолка. Находки превращаются в проверки. \sImpl
  \item[Выпуск без простоя.] Выпуск под трафиком не теряет ни одного запроса; пересоздание кромки
    и базы данных есть операции с окном, измеряемым при каждой отправке против потолков
    (0,3~секунды для кромки на одном хосте, 0,6~секунды для падения базы данных), а оставшаяся
    половина задачи о кромке открыто числится в книге готовности. \sImpl
\end{description}

\subsection{Нагрузка: что измерено, что спроецировано и на чём держится число}

В репозитории встречаются числа трёх родов, и документ держит их порознь. \emph{Измеренные}
получены на названной машине, ноутбуке с восьмиядерным процессором Эппл, и несут отметку версии.
\emph{Смоделированные} получены на стенде национального моделирования: синтетическая страна,
прогнанная через настоящие процедуры в указанном масштабе. \emph{Спроецированные} есть арифметика
над измеренными, и сам репозиторий так их и помечает, называя допущение. Ни один промышленный
кластер не измерялся.

{\footnotesize
\setlength{\tabcolsep}{4pt}\renewcommand{\arraystretch}{1.12}
\rowcolors{2}{tabstripe}{white}
\begin{xltabular}{\textwidth}{>{\RaggedRight\arraybackslash\hspace{0pt}}X >{\RaggedRight\arraybackslash\hspace{0pt}}p{1.8in} >{\RaggedRight\arraybackslash\hspace{0pt}}p{0.8in} >{\RaggedRight\arraybackslash\hspace{0pt}}p{1.15in}}
\caption{Показатели производительности на тех версиях, где они измерены. Каждая строка относится к одному узлу на эталонном ноутбуке, если не помечена как моделирование или проекция. Скорость приёма событий и скорость криптографической проверки различаются примерно в 35~раз, и собственная история репозитория однажды этот разрыв скрыла (раздел~\ref{sec:senses}). Строки, которые репозиторий не перемерял со времён версии 2, несут свою прежнюю отметку.}\label{tab:perf}\\
\toprule
\rowcolor{tabheader}\thead{Величина} & \thead{Значение} & \thead{Род} & \thead{Источник} \\
\midrule
\endfirsthead
\toprule
\rowcolor{tabheader}\thead{Величина} & \thead{Значение} & \thead{Род} & \thead{Источник} \\
\midrule
\endhead
\midrule\multicolumn{4}{r}{\itshape продолжение на следующей странице}\\
\endfoot
\bottomrule
\endlastfoot
Подпись МЛ-ДСА-65, одно ядро & около 2\,110 в секунду & измерено & стенд производительности, \code{в9.278} \\
Проверка МЛ-ДСА-65, один свидетель, одно ядро & около 7\,840 в секунду & измерено & стенд производительности, \code{в9.278} \\
Проверка МЛ-ДСА-65, два свидетеля, одно ядро & около 740 в секунду & измерено & стенд производительности, \code{в9.278} \\
Доказательство принадлежности, глубина 14 & около 35~мс на изготовление, 13~мс на проверку, 76~КБ & измерено & стенд и книга обоснованности \\
Закрытие эпохи на глубине 24, миллион участников & 2,1~с, десятки мегабайт (было 10,8~с и 2,9~ГБ на тысячу) & измерено & темп эпох, \code{в9.357} \\
Переподписание населения под МЛ-ДСА-87, один исполнитель & около 345 удостоверений в секунду, из них 91~процент времени занимает подпись & измерено & учение о квантовом событии, \code{в9.365} \\
Выдача через приложение & 40 в секунду устойчиво, 95-й процентиль 28~мс & измерено & базовая линия, \code{в9.191} \\
Событие проверки через приложение & 80 в секунду устойчиво, 95-й процентиль 18,5~мс & измерено & базовая линия, \code{в9.191} \\
Задержка проверки на топологии из двух участников при переключении & медиана при примерно 2 запросах в секунду; 95-й и 99-й процентили не приводятся, выборка их не выдерживает & измерено & учение о переключении, \code{в9.387} \\
Регистрация с настоящей подписью, массовый конвейер & около 372 токенов в секунду & моделирование, 1:10\,000 & национальное моделирование, \code{в9.257} \\
Приём событий проверки (запись аудита, без проверки подписи) & около 25\,970 в секунду & моделирование, 1:10\,000 & национальное моделирование \\
Сводки Атласа, от 8\,010 до 80\,010 событий & всякая ограниченная сводка росла медленнее своих данных (от 0,72 до 0,89 от линейного роста) & измерено & замер, \code{в9.402} \\
Проверка одним свидетелем на восьмиядерном узле & около 62\,800 в секунду & проекция, в допущении распараллеливания & арифметика над замером \\
\end{xltabular}
}

Плановые цели дорожной карты (десять миллионов человек для штата; 350~миллионов, пять тысяч
устойчивых и пятьдесят тысяч пиковых проверок в секунду для страны) проверены на \code{в9.384}, и
записаны обе половины ответа. Каждая цель по пропускной способности перекрыта больше чем на
порядок: пику нужно примерно шесть с половиной ядер проверки одним свидетелем, устойчивому темпу
меньше одного, а всплеск в двести тысяч регистраций за день есть девять минут работы одного
подписанта. И вместе с тем система не проработала бы и недели на устойчивом темпе, потому что
идентификатор события проверки был 32-битной последовательностью: два миллиарда строк при пяти
тысячах в секунду есть пять дней, а на пике двенадцать часов, и когда последовательность
исчерпана, ничто не замедляется, просто каждая вставка на этом пути падает. Пять суррогатных
идентификаторов теперь 64-битные, проверка пересчитывает эту арифметику из живой схемы при
каждой отправке, а правило, которое репозиторий из этого вывел, таково: вопрос о ёмкости, на
который ответили в ядрах и секундах, есть лёгкая его половина. Модель также отказывается называть
число ядер \emph{измеренным}: измерена скорость на одно ядро, а умножение ни разу не запускали, и
начиная с \code{в9.387} каждая экстраполяция называет своё допущение. Отдельная модель стоимости
находит, что пропускная способность проверки не есть статья расходов ни при каком реалистичном
национальном масштабе, поскольку сто миллионов человек, проверяемых двенадцать раз в год, это
в среднем тридцать восемь проверок в секунду, и она указывает, какие из её входов суть факты о
коде, а какие догадки о развёртывании. Национальная нагрузка смоделирована в масштабе и прогнана
через настоящие ограничения; в действительности она не запускалась. \sExp

\subsection{Злоупотребления, секреты, цепочка поставок и решения, оставляющие запись}

Квоты для каждого ведомства на выдачу, отзыв и проверку суть необязательные строки, которые
закрепляет триггер на каждом пути записи; при одновременной работе они точны благодаря
рекомендательной блокировке, при превышении громко отказывают на каждой поверхности и привязаны к
идентификатору ведомства, но никогда к человеку; начиная с \code{в9.424} квота есть решение с
причиной, переживающее собственное изменение, а не настройка, которую можно перезаписать.
Оповещения о скорости сравнивают каждое ведомство с его собственной прошлой неделей. Ограничитель
частоты по адресу ИП делит одну корзину в Редисе между всеми рабочими процессами. Операторы после
срока регистрации опознаются по ВебАутн как по второму фактору, сессии регистрируются на сервере с
ограничением по ролям и проверкой источника, и для каждой роли действует своя сетевая политика.
Создание операторской учётной записи до \code{в9.422} не оставляло никакой записи о том, что оно
произошло, а у шести из двадцати одного типа событий аудита, названных схемой, во всём дереве не
было ни одного писателя; теперь писатель есть у каждого, это закреплено проверкой, а первая
версия той проверки ошибалась насчёт трёх из них и об этом говорит. Секреты монтируются файлами,
запечатываются и сменяются по одному жизненному циклу; в промышленном режиме приложение
отказывается запускаться с ключом по умолчанию, а журнал приложения не записывает ни тела форм,
ни значения куки, ни значения токенов, и это было задокументированным обещанием, державшимся
одной дисциплиной, пока проверка не начала читать сами вызовы журналирования. Каждый выпуск несёт
перечни состава с подписанным происхождением, а поиск известных уязвимостей в зависимостях и в
каждом образе собственной сборки запирает сборку. \sImpl

\subsection{Постоянный свидетель и смотрит ли он}

При каждой отправке выполняется двадцать задач НИ, и у большинства есть комментарий, называющий
тот прошлый отказ, ради предотвращения которого она существует. Они прогоняют набор тестов
продукта и слой инвариантов; водят консоль Атласа в безголовом браузере; собирают и поднимают
образы; собирают кромку и доказывают постквантовое рукопожатие против настоящего сертификата;
доставляют оповещение о принуждении через настоящий диспетчер оповещений до обратного вызова;
проверяют передачу трассировки и панели наблюдения; убивают и восстанавливают ведущий узел;
ставят профиль Линукс на два дистрибутива; подписывают внутри токена через программный модуль
ПККС\#11 при двух свидетелях; выкатывают выпуск под трафиком; эвакуируют регион; роняют ведущий
узел, хранилище аренды и участника ЕТЦД в профиле отказоустойчивости под живым потоком записей;
поднимают профиль Кубернетеса и переключают его базу данных; подписывают и проверяют настоящим
МЛ-ДСА-65; связывают федерацией два экземпляра поверх ХТТП; ставят \code{полярис-проверка} так,
как поставил бы посторонний; проверяют типы, тестируют и сверяют на соответствие КСР на
ТайпСкрипте; ищут уязвимости в зависимостях и каждом образе; и поднимают весь промышленный стек
от края до края. \sImpl

Версия 3 добавляет то, чего не могла версия 2: у свидетеля спросили, смотрит ли он. Нашлось
двадцать два учения, которые печатали весь абзац своего вердикта, не записав ни одного случая,
поэтому учение теперь отказывается сообщать гарантию, которую не проверило, а проверка требует от
каждого учения считать свои случаи. Учение, подключённое к задаче НИ, где не было нужной ему среды
выполнения, отказалось давать отчёт, а не объявило непроверяемое дерево чистым, и это было
правильно; но задача после этого падала на собственной подготовке, а не на находке, а именно так
настоящую находку и пропускают сквозь пальцы. Теперь каждая задача ставит то, что запускает, и
проверка это требует. Местный запор, который сопровождающий прогоняет перед отправкой, выполнял
четыре набора тестов, тогда как НИ выполняла восемнадцать, и дважды выпустил сломанное изменение
прежде, чем НИ об этом сказала; теперь он выполняет то же, что НИ, и называет проверку, которая
нужна изменению, по тем путям и обработчикам маршрутов, которые сдвинулись. \sImpl

\subsection{День, когда заканчивается пилот, и день, когда старое удостоверение перестают принимать}

Две операции записаны не абзацами, а процедурами, которые выполняются, потому что каждая из них
есть то обещание, на которое учреждение в действительности соглашается. Настоящее обещание пилота
состоит не в том, что он сработает, а в том, что его можно свернуть, и это обещание отказывает
тихо: стирание превращается в абзац в форме согласия, который никто не исполняет. Сворачивание
(\code{в9.376}) отзывает каждое удостоверение пилота и заменяет каждое имя бессмысленной меткой под
записью о стирании, работающей только на добавление; текст согласия, который видит участник,
порождается из того, что делает код, а не пишется из того, что звучало бы успокоительно; и вся
процедура выполняется против настоящей базы данных при каждой отправке. Вместе с ней закрыт и
пакет документов пилота, который учреждение обязано иметь до начала. \sImpl{} для процедуры;
\sExt{} для пилота, поскольку ни один не проводился.

Новое национальное удостоверение появляется рядом с водительскими правами и паспортом и долгие
годы остаётся второй вещью, которую человек носит с собой. Опасный момент этого периода не
переключение, а закат, день, когда старое удостоверение перестают принимать, потому что в этот
момент новое становится обязательным, хотя никто не решал сделать его таковым. План
сосуществования (\code{в9.381}) обращается с закатом как с решением, под которым лежит нижний
предел, а не как с этапом, который наступает сам, и его вердикт отказывается вычисляться из того,
что видит эмитент: принимают ли доверяющие стороны удостоверение и существует ли пригодный
обходной путь, суть заверения, которые делает оператор, и система их записывает, а не
угадывает. \sImpl

Рядом с этим стоят два сопоставления. Сопоставление с НИСТ СП 800-63 излагает строка за строкой,
на чём держалось бы утверждение об уровне достоверности, причём каждая ссылка разрешается и
прогоняется учением, так что строка не может ссылаться на свидетельство, которого нет; это прямо
не есть утверждение о соответствии, и никакой оценки не проводилось. А доступность для людей с
ограничениями проверяется при каждой отправке в той трети, которую покрывает автоматика, и прямо
названа несделанной в остальном, потому что идентификационная система, которой человек не может
воспользоваться, исключает его из самой идентичности. \sImpl{} для механизмов; \sExt{} для
оценки.

\begin{layercard}{Карточка слоя: эксплуатация}
\qa{Что это}{Стек из пяти служб на четырёх основаниях; отказоустойчивость, резервный регион, реплики, разбиение на разделы, резервное копирование и аварийное восстановление, хаос, выпуски без простоя; защита от злоупотреблений, секреты, запоры цепочки поставок; двадцать задач НИ и два учения по расписанию; сворачивание пилота и план сосуществования.}
\qa{Зачем оно существует}{Правильная система, которая не переживает потерянного диска, потерянного региона, неудачного выпуска или враждебного клиента, не есть инфраструктура, а система, которую нельзя свернуть, не та, с которой учреждению следует начинать.}
\qa{Чему доверяет}{Одноразовым хостам учений как нижней оценке промышленного поведения; тому, как оператор разместил хосты и регионы, и выбранному им способу хранения ключей.}
\qa{Чему отказывает}{Утверждению без учения; числу без отметки; числу ядер, названному измеренным; учению, сообщающему гарантию, которую оно не проверило; контейнеру, работающему от суперпользователя; незакреплённому образу; исправимой критической уязвимости.}
\qa{Что видит}{Метрики, трассы и оповещения о системе. Идентификатор корреляции на каждом запросе никогда не касается аудита как самой записи.}
\qa{Чего не видит}{Человека: квоты, оповещения, метрики и журналы привязаны к ведомству, маршруту и рабочему процессу.}
\qa{Если откажет}{Учение, превысившее свой потолок, роняет сборку; находка хаоса или восстановления становится проверкой; эвакуация региона, выдумавшая или пропустившая строку, падает; промышленный сбой есть то, к чему готовят девять решений оператора из книги готовности.}
\qa{Кем проверяется}{\code{проверка\_автоматика\_отказоустойчивости}, \code{проверка\_восстановление\_между\_регионами}, \code{проверка\_маршрутизация\_реплик}, \code{проверка\_разбиение\_таблиц\_событий}, \code{проверка\_учение\_восстановления\_по\_расписанию}, \code{проверка\_программа\_хаоса}, \code{проверка\_выпуск\_без\_простоя}, \code{проверка\_модель\_ёмкости}, \code{проверка\_защита\_от\_злоупотреблений}, \code{проверка\_у\_каждого\_события\_аудита\_есть\_писатель}, \code{проверка\_журналы\_без\_личных\_данных}, \code{проверка\_укрепление\_контейнеров}, \code{проверка\_поиск\_уязвимостей}, \code{проверка\_учения\_считают\_случаи}, \code{проверка\_план\_сосуществования}.}
\qa{Сейчас или потом}{Сейчас, в учениях. Аппаратный модуль, промышленное железо и настоящий пилот потом.}
\qa{Внутри и снаружи}{Внутри: всё вышеописанное. Снаружи: взгляд системы на саму себя, который нужен оператору без взгляда на людей.}
\end{layercard}

\nextlayer{Оператору, который всем этим управляет, нужно видеть, что делает система, и понимать,
откуда берутся полномочия, и уметь и то и другое так, чтобы система не превратилась в карту своих
людей. Следующий слой есть то, как Полярис наблюдает и понимает сам себя.}
